\SourceLeaf{227}
ein sphärisches Dreieck. Im zweiten Fall liegen sie in derselben Reihenfolge wie die Punkte \((1),(2),(3)\), im dritten in umgekehrter Reihenfolge. Bezeichnen wir die Winkel dieses Dreiecks einfach mit \(L,L',L''\) und das Lot, das auf der Kugeloberfläche vom Punkt \(L''\) zur Seite \(LL'\) gezogen wird, mit \(p\), so gilt
\[
\sin p=\sin L\mathord{.}\sin LL''=\sin L'\mathord{.}\sin L'L'',\quad\text{und}\quad\lambda L''=90^0\mp p
\]
\noindent wobei das obere Vorzeichen für den zweiten Fall gilt, das untere für den dritten. Daraus schließen wir
\[
\begin{aligned}
\pm\Delta&=\sin L\mathord{.}\sin LL'\mathord{.}\sin LL''=\sin L'\mathord{.}\sin LL'\mathord{.}\sin L'L''\\
&=\sin L''\mathord{.}\sin LL''\mathord{.}\sin L'L''
\end{aligned}
\]
\noindent Offenbar lässt sich der erste Fall als im zweiten oder dritten enthalten auffassen. Man sieht leicht, dass \(\pm\Delta\) das Sechsfache des Volumens der Pyramide angibt, die von den Punkten \(L,L',L''\) und dem Kugelmittelpunkt gebildet wird. Schließlich folgt daraus ohne Weiteres, dass derselbe Ausdruck \(\pm\tfrac{1}{6}\Delta\) allgemein das Volumen jeder Pyramide angibt, die vom Koordinatenursprung und von den Punkten mit den Koordinaten \(x,y,z;\ x',y',z';\ x'',y'',z''\) gebildet wird.

\Article{3}
\indent Einer gekrümmten Fläche schreibt man in einem auf ihr liegenden Punkt \(A\) stetige Krümmung zu, wenn die Richtungen aller Geraden von \(A\) zu sämtlichen Punkten der Fläche, die von \(A\) unendlich wenig entfernt sind, unendlich wenig von ein und derselben durch \(A\) verlaufenden Ebene abweichen. Man sagt, dass diese Ebene die gekrümmte Fläche im Punkt \(A\) \textit{berührt}. Lässt sich diese Bedingung in einem Punkt nicht erfüllen, so ist dort die Stetigkeit der Krümmung unterbrochen, wie etwa an der Spitze eines Kegels. Die vorliegenden Untersuchungen beschränken sich auf solche gekrümmten Flächen oder Teile von Flächen, bei denen die Stetigkeit der Krümmung nirgends unterbrochen ist. Hier bemerken wir lediglich, dass die Verfahren zur Bestimmung der Lage der Tangentialebene an singulären Punkten, an denen die Stetigkeit der Krümmung unterbrochen ist, ihre Wirksamkeit verlieren und zu unbestimmten Ergebnissen führen müssen.

\Article{4}
\indent Die Lage der Tangentialebene lässt sich am einfachsten aus der Richtung der Geraden bestimmen, die im Punkt \(A\) senkrecht auf der Ebene steht. Diese Gerade heißt auch Normale der gekrümmten Fläche selbst. Die Richtung dieser Normalen stellen wir durch den Punkt \(L\) auf der Hilfskugel dar und setzen
